\clearpage\begingroup
\emergencystretch=4em\relax\NoSourceHead\SourceLeaf{492b}
\makeatletter\tagsleft@true\let\veqno\@@leqno\makeatother
\renewcommand{\thefootnote}{\arabic{footnote})}\setcounter{footnote}{47}
\newcommand{\BolzaVerticalNotEqual}{\mathrel{\vcenter{\hbox{\begin{tikzpicture}[x=1em,y=1ex,line width=.05em]
\draw(0,.18)--(.68,.18);\draw(0,-.18)--(.68,-.18);\draw(.34,-.6)--(.34,.6);
\end{tikzpicture}}}}}
\ArithmeticHead{35}{R}{GAUSS’ PRINCIPIA GENERALIA THEORIAE FIGURAE FLUIDORUM.}
\begin{center}\textls[100]{C) Die partielle Differentialgleichung für die Flüssigkeitsoberfläche und die Randbedingung.}\end{center}
\noindent\hspace*{\parindent}18. Nachdem so die erste Variation von $W$ und $s$ auf ihre einfachste
Form gebracht war, handelte es sich nunmehr darum, daraus die Bedingungen
erster Ordnung für ein Minimum des Ausdrucks $W$ bei gegebenem Werte von
$s$ abzuleiten. Hier hätte nun \textsc{Gauss} von der \textsc{Euler}schen Regel Gebrauch
machen können. Diese Regel war zwar von \textsc{Euler}\footnote{\textsc{Euler}, Institutiones calculi integralis, Bd. III (1770), Appendix de calc. var., art. 90, Opera,
Ser. I, vol. 13, S. 409.} und \textsc{Lagrange}\footnote{\textsc{Lagrange}, Misc. Taurinensia 1760/61, III. Abteilung, S. 173ff., Oeuvres, Bd. I, S. 356.} nur für
einfache Integrale abgeleitet, aber doch auch schon auf Doppelintegrale angewandt
worden, und zwar beim Problem des Körpers von kleinster Oberfläche
bei gegebenem Volumen. Aber \textsc{Gauss} hat es auch hier vorgezogen, ganz
unabhängig von seinen Vorgängern seinen eigenen Weg zu gehen. Bei der
prinzipiellen Bedeutung der hier vorliegenden Schlüsse verlohnt es sich, ihm
dabei ins einzelne zu folgen.

\textsc{Gauss} schliesst so (Art. 28): Soll die Fläche $U$ den Ausdruck $W$ bei gegebenem
Werte von $s$ zu einem Minimum machen, so darf $W$ für keine unendlich
kleine Formänderung der Fläche, bei der die Begrenzungskurve auf der
Gefässwand verbleibt, und bei der das Volumen $s$ ungeändert bleibt, »d. h.
für die $\delta s=0$«, einen negativen Zuwachs erfahren. Daher muss in jedem
Punkt der Fläche das Element des Doppelintegrals in dem Ausdruck für $W$,
nämlich
\[
dU\mathbin{.}\delta n\mathbin{.}\left[z+a^2\left(\frac{1}{R}+\frac{1}{R'}\right)\right]
\]
— wo wir zur Abkürzung
\begin{equation*}
\delta e\mathbin{.}\cos(4,5)=\delta n\tag{48}
\end{equation*}
gesetzt haben —, dem entsprechenden Element von $\delta s$, nämlich $dU\mathbin{.}\delta n$ proportional
\par\smallskip\noindent\hfill{\small 5\textsuperscript{*}}\par
\clearpage\SourceLeaf{493}
\ArithmeticHead{36}{L}{BOLZA, GAUSS UND DIE VARIATIONSRECHNUNG.}
\noindent sein, d. h. \textls[100]{die freie Flüssigkeitsoberfläche muss der partiellen Differentialgleichung}
\begin{equation*}
z+a^2\left(\frac{1}{R}+\frac{1}{R'}\right)=\text{Const.}\tag{49}
\end{equation*}
\textls[100]{genügen}. »Denn wenn diese Proportionalität nicht stattfände, so könnte man
offenbar den Wert von $W$ durch eine geeignete Veränderung der Form der
Oberfläche bei unverändert bleibender Begrenzung verkleinern«. Leider führt
\textsc{Gauss} diesen Schluss nicht näher aus. Gemeint ist wohl folgendes\footnote{Fast genau in dieser Weise ist der Schluss durchgeführt in der Dissertation von R. \textsc{Reiff}, \textit{Über
den Einfluss der Capillarkräfte auf die Form der Oberfläche einer bewegten Flüssigkeit}, Tübingen 1879,
S. 8. Dass \textsc{Gauss} etwa diesen Gedankengang im Sinn hatte, wird auch durch den Beweis in Art. 31 der
Arbeit »\textit{Allgemeine Lehrsätze in Beziehung auf die im verkehrten Verhältnisse des Quadrats der Entfernung
wirkenden . . . Kräfte}« (1840), Werke, Bd. V, S. 233 wahrscheinlich gemacht; vgl. unten III. Teil, Nr. 2.}: Angenommen,
die Funktion
\[
M=z+a^2\left(\frac{1}{R}+\frac{1}{R'}\right),
\]
(die stillschweigend als stetig vorausgesetzt wird), wäre nicht konstant auf der
Fläche $U$, so seien 1 und 2 zwei nicht auf der Begrenzung $P$ liegende Punkte
von $U$, in denen $M$ verschiedene Werte $M_1$ und $M_2$ annimmt. Dann ersetze
man die bestimmten Integrale durch die endlichen Summen, deren Grenzen
sie sind, und variiere nun nur zwei Elemente $dU_1$ und $dU_2$ der Fläche $U$, die
die beiden Punkte 1 bezw. 2 enthalten, aber so, dass $\delta s=0$, und lasse die
ganze übrige Fläche, insbesondere auch die Begrenzung ungeändert. Für eine
solche Variation der Fläche ist dann
\[
\delta s=\delta n_1\mathbin{.}dU_1+\delta n_2\mathbin{.}dU_2=0,\quad
\delta W=M_1\delta n_1\mathbin{.}dU_1+M_2\delta n_2\mathbin{.}dU_2,
\]
wenn $\delta n_1$, $\delta n_2$ die Werte von $\delta n$ in den Punkten 1 und 2 bezeichnen. Hieraus
folgt
\[
\delta W=(M_1-M_2)\delta n_1\mathbin{.}dU_1,
\]
und diese Grösse kann durch passende Wahl des Vorzeichens der willkürlich
bleibenden Grösse $\delta n_1$ negativ gemacht werden. Es muss also die Annahme
$M_1\BolzaVerticalNotEqual M_2$ falsch sein, d. h. $M$ muss auf der ganzen Fläche $U$ konstant sein.

Verbindet man den Grundgedanken dieses Beweises mit dem Mittelwertsatz
für Doppelintegrale, so erhält man leicht einen auch heutigen Anforderungen
\clearpage\SourceLeaf{494}
\ArithmeticHead{37}{R}{GAUSS’ PRINCIPIA GENERALIA THEORIAE FIGURAE FLUIDORUM.}
\noindent an Strenge genügenden Beweis. Hiermit nahe verwandt ist der Beweis, den
H. \textsc{Weber}\footnote{\textsc{Ostwalds} Klassiker der exakten Wissenschaften, Nr. 135, S. 70, 71.} in den Anmerkungen zur deutschen Ausgabe der \textsc{Gauss}schen
Arbeit gegeben hat; dabei wird durch einen Kunstgriff die Anwendung des
Mittelwertsatzes umgangen.

Derselbe Gedanke liegt auch den beiden Beweisen für das Fundamentallemma
für isoperimetrische Probleme zugrunde, die J. \textsc{Bertrand} (1842) und
P. \textsc{du Bois Reymond} (1879) gegeben haben\footnote{Vgl. \textit{Encyklopädie der mathematischen Wissenschaften}, Bd. II 1 (\textsc{Kneser}), S. 581. Es handelt
sich um den ersten Beweis von \textsc{du Bois Reymond}, dessen Grundgedanken \textsc{du Bois Reymond} selbst auf
die oben genannte Dissertation von \textsc{Reiff} zurückführt, sodass hier eine mittelbare Einwirkung von \textsc{Gauss}
festzustellen sein dürfte.}, sowie dem Beweis der \textsc{Euler}schen
Regel bei Doppelintegralen, den A. \textsc{Kneser} in § 64 seines \textit{Lehrbuchs der Variationsrechnung}
(Braunschweig 1900) gegeben hat.

19. Es wird jetzt vorausgesetzt (Art. 29), dass die Fläche $U$ die als notwendig
nachgewiesene partielle Differentialgleichung (49) erfüllt. Dann reduziert
sich die Variation von $W$ auf
\begin{equation*}
\delta W=-\int dP\mathbin{.}\delta\nu\mathbin{.}(a^2\cos i-a^2+2\beta^2),\tag{50}
\end{equation*}
wo wir zur Abkürzung: $\delta e\mathbin{.}\cos(5,8)=\delta\nu$ gesetzt haben. Falls $\beta^2\EqualAboveLess a^2$, geht
der Ausdruck für $\delta W$ durch Einführung des durch die Bedingungen.
\begin{equation*}
\sin\frac{A}{2}=\frac{\beta}{a},\quad 0<A<\pi\tag{51}
\end{equation*}
definierten Winkels $A$ über in
\[
\delta W=a^2\int dP\mathbin{.}\delta\nu\mathbin{.}(\cos A-\cos i).
\]

Nun ist nach der Bedeutung der Richtungen 5 und 8 die Funktion $\delta\nu$
an einer Stelle der Begrenzung $P$ positiv oder negativ, jenachdem bei der
betrachteten virtuellen Bewegung die Flüssigkeit an der betreffenden Stelle
über die Begrenzung $P$ hinaustritt oder von ihr zurückweicht. Hieraus schliesst
\textsc{Gauss}, dass entlang der ganzen Begrenzungskurve $P$ die Randbedingung
\begin{equation*}
i=A\tag{52}
\end{equation*}
erfüllt sein muss. Denn wäre in einem Stück $L$ der Begrenzung $i<A$
(bezw. $i>A$), so könnte man durch eine virtuelle Bewegung der Flüssigkeit,
\clearpage\SourceLeaf{495}
\ArithmeticHead{38}{L}{BOLZA, GAUSS UND DIE VARIATIONSRECHNUNG.}
\noindent welche die ganze Begrenzungskurve $P$ mit Ausnahme des Stückes $L$ ungeändert
lässt, für dieses Stück aber von der ersten (bezw. von der zweiten) Art
ist, $\delta W$ negativ machen.

Es ist interessant, dass \textsc{Gauss} auch hier sich nicht mit dem landläufigen
Schluss: $\delta W$ muss verschwinden, also muss wegen der Willkürlichkeit von $\delta\nu$
der Faktor von $\delta\nu$ identisch verschwinden, begnügt, sondern einen Beweis für
notwendig hält, und zwar ist der Grundgedanke dieses Beweises derselbe, der
dem \textsc{du Bois Reymond}schen Beweis\footnote{Mathematische Annalen, Bd. XV (1879), S. 297, 300.} des Fundamentallemmas für einfache
Integrale zugrunde liegt.

Wenn $\beta^2>a^2$, so ist der Faktor von $\delta\nu$ unter dem Integral (50) stets
positiv und daher kann $\delta W$ stets durch eine virtuelle Bewegung erster Art
negativ gemacht werden. In diesem Fall kann also ein Minimum des Integrals
$W$ nicht stattfinden.

\textls[100]{Mit der Gleichung} $i=A$ \textls[100]{war zum ersten Mal in der Geschichte der Variationsrechnung für ein zweidimensionales Variationsproblem\linebreak die Randbedingung bewiesen.}
Zugleich dürfen wir
als ein Nebenresultat feststellen, dass \textsc{Gauss} die Beweise für das Fundamentallemma
der Variationsrechnung und für die \textsc{Euler}sche Regel auf eine höhere
Stufe der Strenge gebracht hat, als vor ihm erreicht worden war. —
\par\clearpage\endgroup
